% Main story: motivation -> question -> pipeline -> evidence -> decision.
\begin{frame}[plain]
\begin{tikzpicture}[remember picture,overlay]\foot
\at{1}{1}{18}{\fontsize{15}{18}\selectfont\color{teal}COMPUTATIONAL MATERIALS DESIGN}
\at{1}{3.2}{18}{\fontsize{43}{49}\selectfont\bfseries From redox molecules\\[.2cm]to grafted polymers}
\at{1}{7.6}{16.5}{\fontsize{26}{33}\selectfont Which molecular advantages survive attachment and nearby redox groups?}
\fig{18}{3.1}{13}{9.2}{../results/figures/slides/2026-10-30_tempo/pmdi_neighbor_geometry.pdf}
\at{19}{12.1}{11}{\fontsize{17}{21}\selectfont Actual computed PMDI-grafted\\open-chain model}
\claim{Select a useful material, not just a promising molecule.}
\end{tikzpicture}
\end{frame}

\begin{frame}[plain]
\begin{tikzpicture}[remember picture,overlay]\foot
\head{Motivation}{Larger redox carriers can aid separation}
\at{1.2}{3.5}{12}{\fontsize{26}{32}\selectfont\bfseries Store charge on redox motifs.\\[.6cm]Use a larger carrier to aid separation.}
\at{1.2}{8.6}{12}{\fontsize{22}{28}\selectfont Route: attach motifs\\to a preformed polymer\\with benzyl chloride sites.}
\begin{scope}[shift={([xshift=15cm,yshift=-6.5cm]current page.north west)}]
\draw[ink,line width=2pt,rounded corners] (0,0) rectangle (5,-5);
\draw[ink,line width=2pt,rounded corners] (9,0) rectangle (14,-5);
\fill[teal!10] (.1,-1) rectangle (4.9,-4.9);\fill[gold!10] (9.1,-1) rectangle (13.9,-4.9);
\draw[ink,line width=2pt] (6,-.4) rectangle (8,-4.6);
\foreach \y in {1,2,3,4}{\draw[ink!35,line width=1pt] (6,-\y)--(8,-\y);}
\foreach \x/\y in {1/-2,3/-3.5,2/-4}{\fill[teal!30] (\x,\y) circle(.48);\foreach \a in {0,72,...,288}{\fill[teal] (\x,\y) ++(\a:.4) circle(.09);}}
\foreach \x/\y in {10/-2,12/-3.5,11/-4}{\fill[gold!30] (\x,\y) circle(.48);\foreach \a in {0,72,...,288}{\fill[gold] (\x,\y) ++(\a:.4) circle(.09);}}
\node[font=\fontsize{19}{23}\selectfont,align=center] at (7,-6.3) {Redox carriers\quad $|$\quad size-selective separator};
\node[font=\fontsize{14}{18}\selectfont] at (7,1) {Conceptual flow-battery architecture};
\end{scope}
\claim{The design spans redox chemistry and polymer context.}
\src{Context: Janoschka et al., Nature 2015, doi:10.1038/nature15746; Montoto et al., JACS 2016, doi:10.1021/jacs.6b06365. Cartoon is illustrative.}
\end{tikzpicture}
\end{frame}

\begin{frame}[plain]
\begin{tikzpicture}[remember picture,overlay]\foot
\head{The scientific question}{A motif must remain useful after grafting}
\foreach \x/\num/\label in {1/01/Useful redox properties,11.3/02/Feasible attachment,21.6/03/Weak motif interactions}{
\fill[mist,rounded corners=8pt] ([xshift=\x cm,yshift=-4cm]current page.north west) rectangle ++(9.4,-8.3);
\at{\x}{4.4}{9}{\fontsize{42}{46}\selectfont\color{teal}\num}
\at{\x}{6.3}{9}{\fontsize{25}{31}\selectfont\bfseries\label}}
\at{1.3}{8.4}{8.5}{Voltage window\\Accessible capacity\\Stable charge states}
\at{11.6}{8.4}{8.5}{Exact precursor\\Defined reaction conditions\\Core retained after grafting}
\at{21.9}{8.4}{8.5}{Spacing and conformation\\Motif contact strength\\Intersite coupling}
\claim{A candidate must pass every required axis.}
\end{tikzpicture}
\end{frame}

\begin{frame}[plain]
\begin{tikzpicture}[remember picture,overlay]\foot
\head{Pipeline}{Route screening links UMA starts to solution DFT}
\foreach \x/\w/\label in {1/6/Exact models,8.5/6/Route review,16/6/UMA,23.5/7.3/DFT (MeCN)}{
\fill[mist,rounded corners=6pt] ([xshift=\x cm,yshift=-4cm]current page.north west) rectangle ++(\w,-4.2);
\at{\x}{4.4}{\w}{\fontsize{24}{29}\selectfont\bfseries\label}}
\at{1}{6}{5.9}{\fontsize{19}{24}\selectfont Exact graphs\\Precursor/product\\Oligomer occupancy}
\at{8.5}{6}{5.9}{\fontsize{19}{24}\selectfont Route evidence\\Defined conditions\\Screen before QM}
\at{16}{6}{5.9}{\fontsize{19}{24}\selectfont Gas-phase ranking\\Seeds to relaxed poses\\Charge and spin}
\at{23.5}{6}{7.2}{\fontsize{19}{24}\selectfont Optimize geometry\\Electronic energies\\Hessian + qRRHO}
\foreach \x in {7.3,14.8,22.3}{\draw[-{Latex[length=4mm]},teal,line width=2pt] ([xshift=\x cm,yshift=-6.1cm]current page.north west)--++(.85,0);}
\draw[teal,line width=1.4pt] ([xshift=27cm,yshift=-8.2cm]current page.north west)--++(0,-1.25)--++(-21.5,0);
\foreach \x in {5.5,16,27}{\draw[-{Latex[length=3mm]},teal,line width=1.4pt] ([xshift=\x cm,yshift=-9.45cm]current page.north west)--++(0,-.85);}
\foreach \x/\label in {1/Molecular response,11.5/Grafting $\Delta G$,22/Neighbor effects}{
\fill[mist,rounded corners=6pt] ([xshift=\x cm,yshift=-10.4cm]current page.north west) rectangle ++(8.8,-3.3);
\at{\x}{10.75}{8.7}{\fontsize{22}{27}\selectfont\bfseries\label}}
\at{1}{12}{8.7}{\fontsize{22}{27}\selectfont $E_1$, $E_2$ vs Fc/Fc$^+$\\[.1cm]\fontsize{18}{22}\selectfont Directional solution response}
\at{11.5}{12}{8.7}{\fontsize{22}{27}\selectfont $\Delta G_{\mathrm{graft}}$\par Defined base and byproducts}
\at{22}{12}{8.7}{\fontsize{22}{27}\selectfont Neighbor $\Delta\Delta G$\par Local interaction $E_{\mathrm{int}}$}
\claim{UMA speeds the search; checked DFT supports the decision.}
\src{UMA is a gas-phase structural accelerator, not the source of the reported solution free energies. Each exact polymer graph is evaluated separately. The tested Nitro force-repeatability gate fails; the standard-inference/V100 profile remains separate.}
\end{tikzpicture}
\end{frame}

\begin{frame}[plain]
\begin{tikzpicture}[remember picture,overlay]\foot
\head{Candidate expansion}{A systematic grid anchors a 904-graph pool}
\fig{1}{3.6}{20}{8.1}{../results/figures/slides/2026-10-30_tempo/library.pdf}
\at{23}{4.1}{7}{\fontsize{49}{54}\selectfont\color{teal}\bfseries 904}
\at{23}{6.2}{7}{exact graphs in the registry}
\at{23}{8.5}{7}{\bfseries Design variables\\[.3cm]\normalfont Core\\Linker\\Cap\\Ring substituent}
\at{1}{12.8}{28}{\fontsize{22}{27}\selectfont The 837-grid subset spans core, linker, cap and ring choices.}
\claim{Enumeration feeds exact-route and quantum checks.}
\src{The systematic subset contains 459 viologens, 135 PMDIs and 243 NDIs. The unified registry has 904 graphs; the baseline viologen now has three checked native states; most designs await quantum work. Comparable BQ/AQ and oxidation-active coverage remains incomplete.}
\end{tikzpicture}
\end{frame}

\begin{frame}[plain]
\begin{tikzpicture}[remember picture,overlay]\foot
\head{Key result 1: molecular design}{PMDI voltage responds to ring substitution}
\fig{1}{3.4}{20}{10.3}{../results/figures/slides/2026-10-30_tempo/matched_pmdi.pdf}
\fig{23}{3.9}{7.5}{5.5}{../results/figures/slides/2026-10-30_tempo/pmdi_model.pdf}
\at{23}{10}{7.5}{\bfseries Same baseline\\[.3cm]\normalfont PMDI core\\Ethyl cap\\Local benzyl graft}
\claim{Matched comparisons identify molecular tuning levers.}
\src{Computed first reductions in MeCN, relative to the exact baseline. Four illustrative contrasts from 18 systematic three-core comparisons; all 26 cohort contrasts reconstruct.}
\end{tikzpicture}
\end{frame}

\begin{frame}[plain]
\begin{tikzpicture}[remember picture,overlay]\foot
\head{Key result 2: validation}{The benchmark error scale is about 0.25 V}
\fig{1.5}{3.2}{16}{11}{../results/figures/slides/2026-10-30_tempo/redox_validation.pdf}
\at{20}{4}{10}{\fontsize{52}{58}\selectfont\color{teal}\bfseries 0.247 V}
\at{20}{6.2}{10}{Unfitted mean absolute error}
\at{20}{8.3}{10}{\fontsize{28}{33}\selectfont\bfseries 39 endpoints}
\at{20}{9.65}{10}{\fontsize{28}{33}\selectfont\bfseries 30 molecules}
\at{20}{11}{10}{\fontsize{28}{33}\selectfont\bfseries 7 studies}
\at{20}{12.3}{10}{\fontsize{19}{24}\selectfont MeCN; RMSE = 0.279 V}
\claim{The benchmark defines a limited molecular accuracy context.}
\src{Retrospective exact-free-molecule panel. Molecule-cluster 95\% MAE interval: 0.213--0.286 V. No NDI endpoint; grafted-polymer accuracy is not established.}
\end{tikzpicture}
\end{frame}

\begin{frame}[plain]
\begin{tikzpicture}[remember picture,overlay]\foot
\head{Exact enumerated example}{One exact viologen links redox and grafting}
\fig{1.2}{4.3}{12.3}{5.4}{../results/figures/candidates/viologen_exact_graft.pdf}
\at{1.5}{10.3}{12}{\fontsize{22}{28}\selectfont Exact capped graft graph}
\at{1.5}{11.5}{12}{\fontsize{22}{28}\selectfont Grafting $\Delta G=\ViologenGraftDG$ kJ/mol}
\at{1.5}{12.8}{12}{\fontsize{20}{25}\selectfont Monomer $\Delta G_{\mathrm{disp}}=+\MolecularDisproportionation$ kJ/mol}
\fig{14.4}{3.5}{16.2}{10.5}{../results/figures/manuscript/2026-10-31_exact_CPU_example/viologen_native_redox.pdf}
\claim{One graph connects enumeration, route, redox and grafting.}
\src{Capped local model, MeCN, 298.15 K; three native states and matched Fc, no fitted offset. Bars show rotor-setting sensitivity. Grafting is a finite local-site separated-ion model; occupancy and coupling remain pending. Drawing shows resting connectivity; no charge/spin localization is inferred.}
\end{tikzpicture}
\end{frame}

\begin{frame}[plain]
\begin{tikzpicture}[remember picture,overlay]\foot
\head{Exact-candidate follow-up}{TEMPO needs a directional response model}
\fig{1.2}{3.5}{11.3}{7.6}{../results/figures/slides/2026-10-30_tempo/tempo_exact_route.pdf}
\at{1.5}{11.6}{10.9}{\fontsize{23}{29}\selectfont Reviewed attachment route}
\at{1.5}{12.6}{10.9}{\fontsize{23}{29}\selectfont $E_{\mathrm{ox}}=+0.579$ V vs Fc}
\fig{13}{4.4}{18}{7.9}{../results/figures/slides/2026-10-30_tempo/tempo_solution_relaxation.pdf}
\at{14.5}{12.2}{16}{\fontsize{23}{29}\selectfont Directional asymmetry: \textbf{\TempoAsymmetryPercent{}\%}\\Above the declared 5\% diagnostic}
\claim{Keep both directions; one symmetric parameter is unqualified.}
\src{Exact capped TEMPO graft; native-SMD frozen-slow electronic response in MeCN at two checked geometry-level minima. Not grafting free energy or a measured rate. No gas-inner or Born-outer term is added.}
\end{tikzpicture}
\end{frame}

\begin{frame}[plain]
\begin{tikzpicture}[remember picture,overlay]\foot
\head{Key result 3: grafting}{The base defines the grafting reaction}
\at{1.4}{3.2}{29}{\fontsize{25}{31}\selectfont Motif--OH + polymer--CH$_2$Cl + B $\longrightarrow$ polymer--CH$_2$--O--motif + BH$^+$ + Cl$^-$}
\fig{2}{5.1}{27}{8}{../results/figures/slides/2026-10-30_tempo/base_bq.pdf}
\claim{Base-required BQs remain candidates for further evaluation.}
\src{Local graft-site calculations, MeCN, 298.15 K. Computed DBU capture is a shared offset of $-102.78$ kJ/mol; rankings are unchanged. pKa validation exposes an approximately 23 kJ/mol discrepancy. Signs remain conditional; selectivity and yield are untested.}
\end{tikzpicture}
\end{frame}

\begin{frame}[plain]
\begin{tikzpicture}[remember picture,overlay]\foot
\head{Beyond the monomer}{Four states isolate the neighbor effect}
\foreach \x/\y/\state/\a/\b in {2/4/00/0/0,17/4/10/1/0,2/9/01/0/1,17/9/11/1/1}{
\begin{scope}[shift={([xshift=\x cm,yshift=-\y cm]current page.north west)}]
\node[anchor=west,font=\fontsize{24}{29}\selectfont\bfseries] at (0,.35) {$G_{\state}$};
\draw[ink,line width=2pt] (0,-1.5)--(1,-1)--(2,-1.5)--(3,-1)--(4,-1.5)--(5,-1)--(6,-1.5)--(7,-1)--(8,-1.5)--(9,-1);
\draw[ink,line width=1.5pt] (3,-1)--(3,-2.4);\draw[ink,line width=1.5pt] (5,-1)--(5,-2.4);
\ifnum\a=1\fill[teal] (3,-2.8) circle(.45);\else\node[font=\fontsize{19}{22}\selectfont] at (3,-2.8) {Cl};\fi
\ifnum\b=1\fill[gold] (5,-2.8) circle(.45);\else\node[font=\fontsize{19}{22}\selectfont] at (5,-2.8) {Cl};\fi
\end{scope}}
\at{1}{13.6}{29}{\centering\fontsize{29}{34}\selectfont $\Delta\Delta G=G_{11}+G_{00}-G_{10}-G_{01}$}
\claim{Measure how a nearby graft changes the next grafting step.}
\src{Schematic of the exact occupancy cycle, not an atomistic drawing. All backbones are open chains. Cl denotes an unreacted benzyl-chloride handle; colored circles denote grafts. Actual loading and spacing are unknown.}
\end{tikzpicture}
\end{frame}

\begin{frame}[plain]
\begin{tikzpicture}[remember picture,overlay]\foot
\head{Key result 4: a linked PMDI example}{A sampled PMDI neighbor has a modest effect}
\fig{1}{3.3}{21}{10.6}{../results/figures/slides/2026-10-30_tempo/pmdi_neighbor_geometry.pdf}
\at{23}{4.1}{8}{\fontsize{46}{52}\selectfont\color{teal}\bfseries $-1.86$}
\at{23}{6.1}{8}{kJ/mol\\Neighbor $\Delta\Delta G$}
\at{23}{9.1}{8}{\bfseries 16.30 \AA\\[.2cm]\normalfont Core separation}
\at{23}{12}{8}{\fontsize{19}{24}\selectfont Four matched free energies; one minimum per graph.}
\claim{The finite PMDI cycle quantifies a polymer-context effect.}
\src{Exact ethyl-PMDI design. Local redox: $E_1=-1.076$, $E_2=-1.977$ V vs Fc. First-site attachment is 5.34/6.11 kJ/mol less favorable than the capped model. Whole-chain nonadditivity includes backbone and solvent response.}
\end{tikzpicture}
\end{frame}

\begin{frame}[plain]
\begin{tikzpicture}[remember picture,overlay]\foot
\head{Key result 5: unwanted association}{A spacer can still permit strong contact}
\fig{1}{3.3}{20}{10.8}{../results/figures/slides/2026-10-30_tempo/local_motif_contacts.pdf}
\at{22}{4.3}{8.5}{\fontsize{35}{42}\selectfont\color{teal}\bfseries $-25$ to $-26$}
\at{22}{6.1}{8.5}{kJ/mol\\Close phthalimide poses}
\at{22}{9.2}{8.5}{\fontsize{35}{42}\selectfont\bfseries $\approx 0$}
\at{22}{11}{8.5}{kJ/mol\\Separated poses}
\claim{A sequence spacer does not guarantee weak contact.}
\src{Frozen capped-fragment electronic interactions in MeCN, $E_{AB}-E_A-E_B$. Whiskers show cap-length sensitivity, not statistical intervals. Phthalimide is a control. No association free energy, population, basis-superposition correction or electronic coupling is established.}
\end{tikzpicture}
\end{frame}

\begin{frame}[plain]
\begin{tikzpicture}[remember picture,overlay]\foot
\head{Candidate selection}{Every required axis must pass}
\fig{1}{3.4}{21}{10.5}{../results/figures/slides/2026-10-30_tempo/descriptor_projection.pdf}
\at{24}{4.2}{7}{\fontsize{50}{56}\selectfont\color{teal}\bfseries 8}
\at{24}{6.2}{7}{Molecular descriptor leaders}
\at{24}{9.1}{7}{\fontsize{50}{56}\selectfont\color{gold}\bfseries 0}
\at{24}{11.1}{7}{Demonstrated all-axis candidates}
\claim{Descriptor leaders still need full material assessment.}
\src{Projection of the three-objective comparison over 56 reduction-active local models. Nominal molecular capacity excludes polymer/loading penalties. Route, stability, accessibility, occupancy and interaction evidence remain incomplete.}
\end{tikzpicture}
\end{frame}

\begin{frame}[plain]
\begin{tikzpicture}[remember picture,overlay]\foot
\head{What we have learned}{Three lessons guide the next designs}
\at{1.3}{4}{2}{\fontsize{39}{43}\selectfont\color{teal}1}
\at{4}{4}{26}{\fontsize{29}{35}\selectfont\bfseries Tune redox chemistry systematically.\\[.2cm]\fontsize{22}{27}\selectfont\normalfont Ring substitutions shift voltage; benchmark error is about 0.25 V.}
\at{1.3}{7.3}{2}{\fontsize{39}{43}\selectfont\color{teal}2}
\at{4}{7.3}{26}{\fontsize{29}{35}\selectfont\bfseries Evaluate a defined attachment reaction.\\[.2cm]\fontsize{22}{27}\selectfont\normalfont Base and speciation affect feasibility; a reference reaction is insufficient.}
\at{1.3}{10.6}{2}{\fontsize{39}{43}\selectfont\color{teal}3}
\at{4}{10.6}{26}{\fontsize{29}{35}\selectfont\bfseries Test the local polymer environment.\\[.2cm]\fontsize{22}{27}\selectfont\normalfont A spacer does not guarantee weak interaction.}
\claim{Next: broader chemistry and complete polymer-context evidence.}
\end{tikzpicture}
\end{frame}
